% Section 4 (writer: steep-b). Source: research/tracks/steeper/notes-final.md §2.1-§2.4 (Setting 2.0, Lemmas 2.1-2.4,
% Thm 2, Cors 2.5, 2.7, Ex 2.7a, Prop 2.8). Notation: ../paper/NOTATION.md (n = samples, N = 2^d, c = lambda/n,
% kappa = machine constant, Phi_1, Phi_2 for St's M_1, M_2, \lev for St's T, epsilon for St's eta).
\section{The teacher's-mistakes theorem}\label{sec:teacher}

H\"anni asks how a learner can grasp what a teacher means rather than become a predictor of what the teacher would say; Solomonoff induction ``would be making the same mistakes the person would make'' \citep{hanni2026notes}. His remedy is the steeper objective of \cref{def:setup:objectives} over product hypotheses: ``if there are explainable mistakes on a simple model, you want the simple model that doesn't predict the mistakes'' \citep{hanni2026notes}. This section makes his example rigorous. \Cref{thm:teacher:main} quantifies over \emph{all} product hypotheses and shows that, for error rules that are incompressible relative to the model, every near-optimal hypothesis is, up to explicit slack, the simple model with noise, on a window of $c=\lambda/n$ of more than three orders of magnitude at his sizes. \Cref{sec:limits} shows where this stops. Proofs are in \cref{app:teacher}.

\subsection{H\"anni's calculation}\label{sec:teacher:hanni}

\begin{remark}[H\"anni's calculation]\status{computed}\src{steeper \S2.1; checks c1, T5 re-run}\label{rem:teacher:hanni}
H\"anni gives two versions of his example \citep{hanni2026notes,hanni2026lesswrong}.
\begin{enumerate}[label=(\arabic*)]
\item \emph{Version 1} (total bits): ``a simple model of $100$ bits which predicts $100$ bits per specification bit, on top of which there is a systematic error model of $100$ bits which predict $10$ bits per specification bit''. Coin, simple model, and simple model plus errors lie on a hull with slopes $100$ and $10$, so the simple model alone is selected exactly for $\lambda\in(10,100)$.
\item \emph{Version 2} (per sample): a $100$-bit ``simple model'' with ``systematic human error'' on ``a $2^{-10}$ fraction of the inputs, with the set of those inputs taking $10000$ bits to specify''. Against the ``model that learns the errors'' ($10{,}100$ bits, loss $0$) and the ``50/50 random distribution'', the ``good model'' (the simple model with flip probability $p=2^{-10}$, $10$ bits for $p$) has rate $(1-H(p))/110=8.989\cdot10^{-3}$ and the errors have rate $H(p)/9990=1.1185\cdot10^{-6}$. So the good model is selected among the three if and only if $c\in(1.12\cdot10^{-6},\,8.99\cdot10^{-3})$, a ratio of about $8000$; plain MDL ($\lambda=1$) switches to the memoriser at $n\approx8.94\cdot10^5$.
\end{enumerate}
A direct computation (code in the accompanying repository) confirms this arithmetic; his ``$\approx p\log_2(1/p)\approx1/100$'' is $0.0098$ against the exact $H(2^{-10})=0.0112$, which does not matter. The calculation compares three hypotheses only. It does not exclude partial models, partial memorisers or mixtures, and H\"anni asks: ``Maybe I'm just failing to construct the right `universal hypothesis' for this example?'' \citep{hanni2026lesswrong}.
\end{remark}

\subsection{The instance}\label{sec:teacher:instance}

\begin{definition}[The teacher's-mistakes instance]\src{steeper Setting 2.0}\label{def:teacher:instance}
Fix integers $1\le m\le d$ and $2\le r\le m$, and put $p:=2^{-r}$ and $s_0:=p2^m$.
\begin{enumerate}[label=(\alph*)]
\item \emph{Model, errors, teacher.} The \emph{simple model} is the parity $f_a(x)=\langle a,x\rangle\bmod2$ on $\{0,1\}^d$, for some $a\in\{0,1\}^d$ with $K(a)\ge d$. The \emph{error rule} is a set $S\subseteq\{0,1\}^m$ with $|S|=s_0$, and the \emph{error set} is $E:=S\times\{0,1\}^{d-m}$, the inputs whose first $m$ bits lie in $S$. The \emph{teacher} is $y^*:=f_a\oplus\ind_E$. For $m<d$ the errors are a systematic rule (a random junta) that generalises to new inputs; for $m=d$ they are an idiosyncratic random set.
\item \emph{Description lengths.} $L_S:=\log\binom{2^m}{s_0}$, $B_S:=K(S\mid a,K(a))$, and the \emph{deficiency} $\delta_S:=L_S-B_S$. The \emph{information} a program $\pi$ holds about the rule beyond the parity is $I(\pi):=B_S-K(S\mid a,\pi)$, so $L_S-K(S\mid a,\pi)=\delta_S+I(\pi)$. One has $2^mH(p)-\log(2^m+1)\le L_S\le2^mH(p)$ \citep{cover2006elements}; $\delta_S\ge-2\log m-2\log s_0-\kappa$, since $S$ can be given by its index among the $s_0$-subsets; and $\Pr_S[\delta_S>j]<2^{-j}$ for a uniformly random rule, since fewer than $2^{L_S-j}$ programs are shorter than $L_S-j$.
\item \emph{Population quantities.} Inputs are uniform on $\{0,1\}^d$, so $\Pr[x\in E]=p$. For a product program $\pi$ and a truncation level $B\in[\log(1/p),\infty]$: $R(\pi):=\Ex_x[-\log\nu_\pi(y^*(x)\mid x)]$, $R_B(\pi):=\Ex_x\min\{B,-\log\nu_\pi(y^*(x)\mid x)\}\le R(\pi)$, $J_c:=c\,\ell+R$, $J^B_c:=c\,\ell+R_B$, $\Phi(c):=\inf_\pi J_c(\pi)$ and $\Phi^B(c):=\inf_\pi J^B_c(\pi)\le\Phi(c)$, the infima over the class. Put $\tau_B:=2^{1-B}$ and $\beta_B:=2^m\log(1+\tau_B)$, both $0$ if $B=\infty$.
\item \emph{Agreement and level.} $u(\pi):=\Ex_x\nu_\pi(f_a(x)\mid x)$. With $\theta_t:=\frac12+2^{-t-1}$, the \emph{level} is $\lev(\pi):=\min\{t\ge1:u(\pi)>\theta_t\}$, and $\lev(\pi):=\infty$ if $u(\pi)\le\frac12$; so $\lev(\pi)=1$ if and only if $u(\pi)>\frac34$. Put $s(t):=2t+2\log(t+1)\le4t$ and $G_j:=1-\log(1+2^{-j}+2p+\tau_B)$ for $j\in[0,\infty]$; $G_j$ is increasing and concave in $j$.
\item \emph{Reference programs} in $\Hcl{\times}$: the \emph{coin} $\pi_{\rm coin}$, with $\ell\le\kappa$ and $R=1$; the \emph{good model} $\pi_g$, which gives $f_a(x)$ probability $1-p$, with $\ell(\pi_g)\le d+\kappa_g$, $\kappa_g:=2\log d+2\log r+\kappa$, and $R(\pi_g)=R_B(\pi_g)=H(p)$; the \emph{memoriser} $\pi_b$, the deterministic $y^*$, with $\ell(\pi_b)\le d+\lceil L_S\rceil+\kappa_b$, $\kappa_b:=2\log d+2\log m+2\log s_0+\kappa$, and $R(\pi_b)=0$.
\item \emph{Slack constant.} $\Lambda_1:=2\log d+2\log m+2\log B\cdot\ind[B<\infty]+\beta_B+\kappa$.
\end{enumerate}
\emph{H\"anni's sizes} are $d=100$, $m=20$, $p=2^{-10}$; then $s_0=1024$ and $L_S=11{,}710.27$ bits, his ``$10000$ bits''.
\end{definition}

The objective is the population objective $J_c$; finite samples are treated in \cref{cor:teacher:sampled}. The deficiency $\delta_S$ enters every slack term below, and every number we give states the $\delta_S$ it assumes; numbers with $\delta_S=0$ concern error rules that are incompressible given the model, and say nothing about compressible ones (\cref{prop:limits:envelope}).

\subsection{Four lemmas}\label{sec:teacher:lemmas}

The first lemma is the reason for choosing parities: a hypothesis that agrees with $f_a$ noticeably more often than a coin determines $a$ up to a few bits.

\begin{lemma}[List decoding for semimeasures]\status{proved; computed}\src{steeper Lemma 2.1 (T5 Lemma 2.5a, generalised); checks c2(b), referee r3}\label{lem:teacher:list}
If $\lev(\pi)=t<\infty$, then $K(a\mid\pi)\le s(t)+2\log d+\kappa$.
\end{lemma}

\emph{Proof idea.} The agreement of $\nu_\pi$ with a parity $f_{a'}$ is $\frac12$ plus half a Fourier coefficient of $\nu_\pi(0\mid x)-\nu_\pi(1\mid x)$, and by Parseval fewer than $2^{2t}$ parities have agreement above $\theta_t$ \citep{goldreich1989hard,kushilevitz1993learning}. The agreement is lower semicomputable, so this list can be enumerated from $\pi$ (\cref{app:teacher:list}).

\begin{lemma}[Risk floor]\status{proved; computed}\src{steeper Lemma 2.2 (T5 Lemma 2.5b); check c2(b)}\label{lem:teacher:floor}
$R_B(\pi)\ge-\log\big(u(\pi)+p+2^{-B}\big)$, with $2^{-\infty}=0$. Hence $\lev(\pi)=t\ge2$ implies $R_B(\pi)\ge G_{t-1}$, and $\lev(\pi)=\infty$ implies $R_B(\pi)\ge G_\infty$.
\end{lemma}

\begin{lemma}[Coding the error rule with a hypothesis]\status{proved}\src{steeper Lemma 2.3; referee r3, q2(c)}\label{lem:teacher:coding}
For integer or infinite $B$, $K(S\mid a,\pi)\le2^mR_B(\pi)+\beta_B+2\log m+2\log B\cdot\ind[B<\infty]+\kappa$.
\end{lemma}

\begin{lemma}[Two-part lower bounds]\status{proved}\src{steeper Lemma 2.4 (T5 Thm 2.5(c))}\label{lem:teacher:twopart}
If $\lev(\pi)=t<\infty$, then
\[
\text{(D1)}\quad\ell(\pi)\ge d-s(t)-\Lambda_1,\qquad\qquad\text{(D2)}\quad\ell(\pi)+2^mR_B(\pi)\ge d+B_S-s(t)-\Lambda_1.
\]
\end{lemma}

\emph{Proof ideas.} \Cref{lem:teacher:floor} is Jensen's inequality; a hypothesis can be right on $E$ for free only on a $p$-fraction of the inputs. \Cref{lem:teacher:coding} turns the per-$z$ average likelihoods into a semimeasure on subsets of $\{0,1\}^m$; the product form of $E$ makes the cost $2^mR_B$, not $2^dR_B$. (D1) is $d\le K(a)\le\ell(\pi)+K(a\mid\pi)+\kappa$; (D2) adds \cref{lem:teacher:coding} and symmetry of information. See \cref{app:teacher:lemmas}, which also proves a sharper floor (\cref{lem:teacher:sharpfloor}).

\subsection{The theorem}\label{sec:teacher:theorem}

Define
\begin{gather*}
\psi_c(R):=c\max\{0,B_S-2^mR\}+R,\qquad \Phi_1(c):=c(d-4-\Lambda_1)+B_S/2^m,\qquad t^*:=(d-\Lambda_1)/4,\\
\Phi_2(c):=\begin{cases}\min\{c(d-8-\Lambda_1)+G_1,\ G_{t^*-1}\}&t^*\ge2,\\ G_1&t^*<2,\end{cases}\\
J_g(c):=c(d+\kappa_g)+H(p)\ge J_c(\pi_g),\qquad J_b(c):=c(d+\lceil L_S\rceil+\kappa_b)\ge J_c(\pi_b),\\
\Delta(c):=\max\{0,\ c(4+\Lambda_1+\kappa_g)+H(p)-B_S/2^m\}.
\end{gather*}
The pieces of $\Phi_2$ are nondecreasing in $c$ and $J_g$ is linear, so $\{c>0:J_g(c)\le\Phi_2(c)\}=(0,c_{\rm hi}]$ with
\[
c_{\rm hi}:=\min\Big\{\frac{G_1-H(p)}{8+\Lambda_1+\kappa_g},\ \frac{G_{t^*-1}-H(p)}{d+\kappa_g}\Big\}\ \ (t^*\ge2),\qquad c_{\rm hi}:=\frac{G_1-H(p)}{d+\kappa_g}\ \ (t^*<2).
\]
$c_{\rm hi}>0$ requires $G_1>H(p)$; with $B=\infty$ this holds if and only if $p<0.0560$.

\begin{theorem}[The teacher's-mistakes theorem]\status{proved}\src{steeper Thm 2}\label{thm:teacher:main}
In the setting of \cref{def:teacher:instance}, with $B\in[\log(1/p),\infty]$, the following hold for every program $\pi$ of $\Hcl{\times}$, and of every product class covered by \cref{prop:related:computable}.
\begin{enumerate}[label=(\roman*),start=0]
\item \emph{Level bounds.} If $\lev(\pi)=1$, then $J^B_c(\pi)\ge c(d-4-\Lambda_1)+\psi_c(R_B(\pi))$. If $\lev(\pi)\ge2$, including $\infty$, then $J^B_c(\pi)\ge\Phi_2(c)$.
\item \emph{Window.} If $c\,2^m>1$ and $c\le c_{\rm hi}$, then $\Phi(c)\ge\Phi^B(c)\ge\min\{\Phi_1(c),\Phi_2(c)\}$ and $J_c(\pi_g)\le\Phi(c)+\Delta(c)$, where $\Delta(c)\le c(4+\Lambda_1+\kappa_g)+(\log(2^m+1)+\delta_S)/2^m$.
\item \emph{Identification.} Let $c\,2^m>1$, $\epsilon\ge0$ and $J_g(c)+\epsilon<\Phi_2(c)$, and let $J^B_c(\pi)\le J_g(c)+\epsilon$; this includes every $\pi$ with $J_c(\pi)\le\Phi(c)+\epsilon$. Then $\lev(\pi)=1$. Put $\Lambda_2:=2^mH(p)-B_S+4+\Lambda_1+\kappa_g\le\log(2^m+1)+\delta_S+4+\Lambda_1+\kappa_g$, $\Lambda_3:=4+2\log d+\kappa$ and
\[
\Delta_\ell:=\kappa_g+\frac\epsilon c+\frac{\Lambda_2+\epsilon/c}{c2^m-1}.
\]
\begin{enumerate}[label=(\alph*)]
\item \emph{It knows the simple model:} $u(\pi)\ge\max\{\frac34,\ 2^{-(J_g(c)+\epsilon)}-p-2^{-B}\}$.
\item \emph{It fits no better than noise:} $R_B(\pi)\ge H(p)-\dfrac{c\Lambda_2+\epsilon}{c2^m-1}$.
\item \emph{It is short:} $\ell(\pi)\le d+\Delta_\ell$.
\item \emph{It holds little information about the error rule:} $I(\pi)\le\Delta_\ell+\Lambda_3$, so $L_S-K(S\mid a,\pi)\le\delta_S+\Delta_\ell+\Lambda_3$.
\item \emph{It does not predict the errors:} for every integer $g\ge1$,
\[
e(\pi):=\Ex\big[\nu_\pi(y^*(x)\mid x)\,\big|\,x\in E\big]\ \le\ 1-u(\pi)+\gamma_g+2^{-g},
\]
where $\gamma_g:=\sqrt{j_g\ln2/(2s_0)}$, $j_g$ is the least integer $j\ge0$ with $j-2\log(j+1)>\Gamma_g$, and $\Gamma_g:=\delta_S+\Delta_\ell+\Lambda_3+1+2\log m+2\log s_0+g+2\log g+\kappa$.
\end{enumerate}
\item \emph{Memoriser regime.} Let $c\,2^m<1$, and let $J^B_c(\pi)\le J_b(c)+\epsilon$ with $J_b(c)+\epsilon<\Phi_2(c)$. Then $\lev(\pi)=1$ and $(1-c2^m)\min\{R_B(\pi),B_S/2^m\}\le\varrho:=c(\lceil L_S\rceil-L_S+\delta_S+4+\Lambda_1+\kappa_b)+\epsilon$. If $\varrho<(1-c2^m)B_S/2^m$, then $R_B(\pi)\le\varrho/(1-c2^m)$ and $e(\pi)\ge2^{-R_B(\pi)/p}-2^{-B}$.
\item \emph{The lower end is sharp.} If $c<H(p)/(\lceil L_S\rceil+\kappa_b+4+\Lambda_1)$, then $J_c(\pi_b)<J_c(\pi)$ for every $\pi$ with $\lev(\pi)=1$ and $R(\pi)\ge H(p)$. Since $L_S=2^mH(p)-O(m)$, this threshold is $2^{-m}\big(1+O((m+\log d+\kappa)/L_S)\big)$.
\end{enumerate}
\end{theorem}

\emph{Proof idea.} For a program of level $1$, (D1) and (D2) give $J^B_c\ge c(d-4-\Lambda_1)+\psi_c(R_B)$, and when $c2^m>1$ the function $\psi_c$ is minimised at $R=B_S/2^m$: each bit of rule information costs $c$ per sample in length and saves only $2^{-m}$ in risk. A program of level $t\ge2$ pays in length through (D1) or in risk through \cref{lem:teacher:floor}; concavity of $G$ reduces the $t$-dependence to two endpoints. Part (ii) reads these inequalities backwards. For (e), if a near-optimal program gave the erroneous labels much probability, $S$ would lie in a small set enumerable from the program (Hoeffding's inequality for sampling without replacement \citep{hoeffding1963probability}), contradicting (d). Parts (iii) and (iv) follow from (D2) and (D1). See \cref{app:teacher:main}.

So in the window $2^{-m}<c\le c_{\rm hi}$, for error rules that are incompressible relative to the model, no product hypothesis of any kind beats the good model by more than $\Delta(c)$, and every near-optimal hypothesis knows the parity, is at most $\Delta_\ell$ bits longer than $d$, holds little information about the rule and gives the erroneous labels little probability. Below $2^{-m}$, by a relative margin, near-optimal hypotheses reproduce the errors. This answers H\"anni's question about a missing ``universal hypothesis'' for such rules.

\begin{remark}[Numbers at H\"anni's sizes]\status{computed}\src{steeper \S2.4 ``Numbers''; checks c3, c6(B), c7(D); referee r1, q4(B)}\label{rem:teacher:numbers}
At $d=100$, $m=20$, $p=2^{-10}$, $B=\infty$ and $\epsilon=0$, a direct computation (code in the accompanying repository) gives \cref{tab:teacher:numbers}. Rows with $\kappa=0$ are illustrations, not results about any particular reference machine; rows with $\kappa=30$ show how the bounds degrade, and every bound is explicit in $\kappa$ and $\delta_S$. The window is $(9.54\cdot10^{-7},\,8.06\cdot10^{-3}]$ at $\kappa=0$ (a ratio of $8453$) and $(9.54\cdot10^{-7},\,3.66\cdot10^{-3}]$ at $\kappa=30$ (ratio $3837$), against H\"anni's three-hypothesis window $(1.12\cdot10^{-6},8.99\cdot10^{-3})$. The memoriser has $e=1$ and the good model $e=2^{-10}$; for $m=24$ ($L_S=187{,}457$) the bound at $c=10^{-4}$ is $e\le0.065$ ($\kappa=0$). Near $2^{-m}$ the memoriser regime is only partly covered: at $\kappa=0$, $\delta_S=0$, (iii) applies only for $c<0.99418\cdot2^{-m}$, and $e\ge\frac12$ needs $c\le0.937\cdot2^{-m}$; in a band of relative width about $6\%$ below $2^{-m}$ nothing useful is proved.
\end{remark}

\begin{table}[t]
\centering\small
\begin{tabular}{@{}cc|ccccc@{}}
\toprule
\multicolumn{7}{@{}l}{(a) Identification at $c=10^{-4}$ (\cref{thm:teacher:main}(ii); $g$ optimised)}\\
\midrule
$\kappa$ & $\delta_S$ & $u\ge$ & $\ell\le$ & $I\le$ & $L_S-K(S\mid a,\pi)\le$ & $e\le$\\
\midrule
0 & 0 & 0.983 & $d+20.4$ & 37.7 & 37.7 & 0.199\\
0 & 20 & 0.983 & $d+20.6$ & 37.9 & 57.9 & 0.218\\
30 & 0 & 0.981 & $d+51.0$ & 98.3 & 98.3 & 0.274\\
30 & 20 & 0.981 & $d+51.2$ & 98.5 & 118.5 & 0.287\\
\bottomrule
\end{tabular}

\medskip
\begin{tabular}{@{}l|cccccc@{}}
\toprule
\multicolumn{7}{@{}l}{(b) Memoriser regime (\cref{thm:teacher:main}(iii), $\delta_S=0$): lower bound on $e$}\\
\midrule
$c\,2^m$ & 0.5 & 0.8 & 0.9 & 0.94 & 0.97 & 0.99\\
\midrule
$\kappa=0$ & 0.955 & 0.831 & 0.658 & 0.483 & 0.223 & 0.010\\
$\kappa=30$ & 0.917 & 0.706 & 0.457 & 0.256 & 0.060 & 0.000\\
\bottomrule
\end{tabular}
\caption{The teacher's-mistakes theorem at H\"anni's sizes $d=100$, $m=20$, $p=2^{-10}$ ($L_S=11{,}710$ bits), with $B=\infty$ and $\epsilon=0$. Rows with $\kappa=0$ are illustrations. Information bounds are in bits, out of $L_S$.}
\label{tab:teacher:numbers}
\end{table}

\begin{remark}[What the theorem relies on]\status{remark}\src{steeper \S2.4 ``The hypotheses Theorem 2 relies on''}\label{rem:teacher:hypotheses}
Without one of the first four of the following hypotheses a step of the proof fails as written; the fifth enters every bound. (1) \emph{A list-decodable model class}: parities admit \cref{lem:teacher:list}, while simple models with short approximations, such as majority-like functions, do not. (2) \emph{The product form $E=S\times\{0,1\}^{d-m}$}: \cref{lem:teacher:coding} averages each $z\in\{0,1\}^m$ over $2^{d-m}$ inputs; for an unstructured $E\subseteq\{0,1\}^d$ it would give $2^dR_B$ bits, useless at $d=100$. (3) \emph{Uniform inputs}: rates, and with them the window, scale with the input law (\cref{prop:limits:inputlaw}). (4) \emph{$p<0.056$}, so that $G_1>H(p)$ (\cref{rem:limits:method}). (5) \emph{A small deficiency $\delta_S$}: every slack term contains it, and \cref{prop:limits:envelope} shows that this is not an artefact of the proof.
\end{remark}

\subsection{Idiosyncratic errors, finite samples and Version 1}\label{sec:teacher:more}

With idiosyncratic errors and the whole truth table, the window becomes a window in $\lambda$ itself. Write $L_E,B_E,\delta_E$ for $L_S,B_S,\delta_S$ when $m=d$.

\begin{corollary}[Idiosyncratic errors, full truth table]\status{proved; computed}\src{steeper Cor 2.5 (T5 Thm 3.5, Cor 3.5a); checks c3, c7(H)}\label{cor:teacher:fulltable}
Take $m=d$, $N=2^d$, $B=\infty$, and let the data be the whole truth table of $y^*$, so $n=N$. A $\lambda$-MAP minimises $\lambda\ell+NR=N\,J_{\lambda/N}$.
\begin{enumerate}[label=(\roman*)]
\item For every $\lambda>1$ with $J_g(\lambda/N)<\Phi_2(\lambda/N)$, in particular every $\lambda\in(1,Nc_{\rm hi})$, every $\lambda$-MAP $\pi$ satisfies \cref{thm:teacher:main}(ii) with $c=\lambda/N$, $\epsilon=0$: $u(\pi)\ge\max\{\frac34,2^{-J_g(\lambda/N)}-p\}$, $NR(\pi)\ge NH(p)-\lambda\Lambda_2/(\lambda-1)$, $\ell(\pi)\le d+\kappa_g+\Lambda_2/(\lambda-1)$, $I(\pi)\le\kappa_g+\Lambda_2/(\lambda-1)+\Lambda_3$, and (e).
\item At $\lambda=1$ the optimum, $J_1(\pi_g)$ and $J_1(\pi_b)$ lie within $O(\log N)+\delta_E$ of each other, so $\lambda=1$ does not decide whether the errors are memorised; lower-order terms and the reference machine do.
\end{enumerate}
At $d=20$, $p=2^{-10}$ ($L_E=11{,}710$ bits) a direct computation gives the window $\lambda\in(1,1.19\cdot10^4]$ at $\kappa=0$ and $(1,6456]$ at $\kappa=30$ (bounds in \cref{app:teacher:fulltable}).
\end{corollary}

The theorem is a population statement. For sampled data, let $n\le2^d$ inputs be drawn uniformly without replacement, so that the data form a partial function, and labelled by $y^*$; let $B<\infty$, $\delta\in(0,1)$, $\epsilon_1:=\log(1/p)\sqrt{\ln(2/\delta)/(2n)}$, $\hat L:=(J_c(\pi_g)+\epsilon_1)/c$ and $\epsilon_n:=\epsilon_1+B\sqrt{(\hat L\ln2+\ln(2/\delta))/(2n)}$.

\begin{corollary}[H\"anni's numbers with sampled data]\status{proved; computed}\src{steeper Cor 2.7, Lemma 2.6; check c3}\label{cor:teacher:sampled}
With probability at least $1-\delta$, every $\lambda$-MAP with $\lambda=cn$ satisfies $J^B_c\le J_g(c)+\epsilon_n$, and hence the conclusions of \cref{thm:teacher:main}(ii) with $\epsilon=\epsilon_n$, whenever $c2^m>1$ and $J_g(c)+\epsilon_n<\Phi_2(c)$.
\end{corollary}

The proof is uniform convergence over all programs with Kraft weights (\cref{lem:teacher:finite}), which is pessimistic: at H\"anni's sizes, with $B=30$, $\delta=0.01$, $\kappa=0$, $\delta_S=0$ and $c=10^{-4}$, the excess is $\epsilon_n/c=9.2$ bits at $n=10^{11}$ and $0.9$ bits at $n=10^{13}$ (\cref{app:teacher:sampled}).

\begin{example}[H\"anni's Version 1, in population form]\status{proved (instance of \cref{thm:teacher:main}); computed}\src{steeper Ex 2.7a; checks c6(C), c6(H)}\label{ex:teacher:version1}
Take $d=100$, $m=10$, $p=2^{-6}$, so $s_0=16$ and $L_S=115.6$ bits. At $n=10^4$ the model saves $n(1-H(p))=8839$ bits for its $100$ bits ($88$ bits per bit) and the error rule saves $nH(p)=1161$ bits for its $116$ bits ($10.0$ bits per bit): Version 1's sizes and rates. The window is $c\in(2^{-10},c_{\rm hi}]$ with $c_{\rm hi}=5.80\cdot10^{-3}$ at $\kappa=0$ and $2.53\cdot10^{-3}$ at $\kappa=30$; at $n=10^4$ and $\kappa=0$ this is $\lambda\in(9.77,58.0]$, against H\"anni's $(10,100)$. Only coarse identification is proved at this size: near-optimal hypotheses have $u>\frac34$, and at $\kappa=0$, $\delta_S=0$, $c=5.5\cdot10^{-3}$ they have $\ell\le d+28.3$ and $I\le45.6$ bits out of $115.6$; near the lower end the information bound exceeds $L_S$, and (e) is vacuous because $s_0=16$ is too small. At $n=10^4$ samples \cref{cor:teacher:sampled} gives nothing (\cref{app:teacher:version1}).
\end{example}

Finally, the power posterior with $\lambda=cn$, not only its maximiser, concentrates on near-optimal programs. Let $\hat J_c(\pi):=c\,\ell(\pi)+\frac1n\sum_{i\le n}-\log\nu_\pi(y_i\mid x_i)$ (\cref{rem:setup:persample}) and $\hat J_g:=\hat J_c(\pi_g)$.

\begin{proposition}[The power posterior concentrates on near-minimisers]\status{proved; computed}\src{steeper Prop 2.8; checks c5(e), c7(F); referee q1(F)}\label{prop:teacher:concentration}
Let $\lambda=cn$ and let $D$ be any data set of size $n$. For every $\epsilon>0$,
\[
\frac{\Pi_\lambda\big(\{\pi:\hat J_c(\pi)\ge\hat J_g+\epsilon\}\,\big|\,D\big)}{\Pi_\lambda(\pi_g\mid D)}\ \le\ 2^{-n\epsilon+(\hat J_g+\epsilon)/c}.
\]
Consequently, for data sampled as in \cref{cor:teacher:sampled}, with probability at least $1-\delta$ all but this fraction of the posterior mass lies on programs that satisfy \cref{thm:teacher:main}(ii) with $\epsilon':=\epsilon+\epsilon_1+B\sqrt{(\hat L'\ln2+\ln(2/\delta))/(2n)}$ in place of $\epsilon$, where $\hat L':=(J_g(c)+\epsilon_1+\epsilon)/c$, provided $J_g(c)+\epsilon'<\Phi_2(c)$.
\end{proposition}

\emph{Proof idea.} $2^{-\lambda\ell(\pi)}\nu_\pi(D)=2^{-n\hat J_c(\pi)}$ and $n\hat J_c(\pi)-\ell(\pi)\ge(1-1/\lambda)\,n\hat J_c(\pi)$; sum over the bad set with Kraft's inequality (\cref{app:teacher:concentration}). With $\hat J_g$ replaced by its population value $0.0232$ ($\kappa=0$), $c=10^{-4}$, $\epsilon=10^{-3}$ and $n=10^6$, the ratio is at most $2^{-758}$.

\paragraph{Open question.}\label{q:teacher:open}
Which simple-model classes give a window? Parities do, through list decoding; majority-like functions have short approximations and presumably do not. Can the window method reach beyond $p<0.056$, towards $p<1/4$ where the kink exists (\cref{prop:algstat:parity}, \cref{rem:limits:method}), with sharper constants? And what can be proved at realistic sample sizes, far below the $n\approx10^{11}$ that uniform convergence needs here?
